\begin{threeparttable}
\begin{tabular}{l cc}
\toprule
\multicolumn{3}{c}{\textit{RDD2: Índice de Focalización de Hogares (IFH) como running variable (Adultos ≥65, siguiendo Bando, Galiani y Gertler 2020)}} \\
\midrule
  & TIENE\_BILLETERA & USA\_BILLETERA \\
\midrule
\addlinespace
\multicolumn{3}{l}{\textit{RDD2\_IFH\_baseline}} \\
  Estimate & 0.0280 & 0.0152 \\
  SE (robust) & (0.0211) & (0.0236) \\
  95\% CI & [-0.0133, 0.0694] & [-0.0310, 0.0614] \\
  N / BW & 5{,}052 / 12.48 & 4{,}479 / 10.76 \\
  Method & rdrobust & rdrobust \\
\addlinespace
\multicolumn{3}{l}{\textit{RDD2\_IFH\_covariates}} \\
  Estimate & 0.0272 & 0.0108 \\
  SE (robust) & (0.0184) & (0.0152) \\
  95\% CI & [-0.0088, 0.0631] & [-0.0189, 0.0405] \\
  N / BW & 4{,}653 / 11.52 & 4{,}505 / 11.03 \\
  Method & rdrobust & rdrobust \\
\bottomrule
\end{tabular}
\begin{tablenotes}
\small
\item \textit{Notes:} La running variable es el Índice de Focalización de Hogares (IFH) de Bernal, Carpio y Klein (2017), construido con pesos de las Tablas A.8/A.9 del Apéndice F y estandarizado en [0--100] por departamento. El umbral varía por departamento según la Tabla A.10 (rango: 33--55). Muestra: adultos mayores (EDAD$\geq$65) en ENAHO 2024 con IFH no nulo (N=13{,}766). Este diseño replica la estrategia de identificación de Bando, Galiani y Gertler (2020) aplicada a la adopción de billetera digital.
\end{tablenotes}
\end{threeparttable}